\section{Homework Week 10}

\begin{exercise}
    Let \(R\) be a ring and let \(V\) be an \(R\)-module. Let \(I\) be an ideal of
    \(R\). Suppose that \(d_I\) is defined on \(V\) and on \(R\). Then:
    \begin{enumerate}
        \item Addition is continuous from \(V\times V\) to \(V\).
        \item Multiplication is continuous from \(R\times V\) to \(V\).
        \item If \(d_I\) is defined on the \(R\)-module \(W\) and
              \(f\in \operatorname{Hom}_R(V,W)\), then \(f\) is continuous.
    \end{enumerate}
    \end{exercise}
    
    \begin{proof}
    Recall that, for an \(R\)-module \(M\), the \(I\)-adic topology induced by
    \(d_I\) has a basis of neighbourhoods of \(x\in M\) given by
    \[
        x+I^nM,\qquad n\geq 1.
    \]
    
    \begin{enumerate}
        \item Let \((v_0,w_0)\in V\times V\). We show that addition is continuous at
              \((v_0,w_0)\). Let
              \[
                  v_0+w_0+I^nV
              \]
              be a neighbourhood of \(v_0+w_0\). If
              \[
                  v-v_0\in I^nV,\qquad w-w_0\in I^nV,
              \]
              then
              \[
                  (v+w)-(v_0+w_0)=(v-v_0)+(w-w_0)\in I^nV.
              \]
              Hence
              \[
                  v+w\in v_0+w_0+I^nV.
              \]
              Therefore addition \(V\times V\to V\) is continuous.
    
        \item Let \((r_0,v_0)\in R\times V\). We show that scalar multiplication is
              continuous at \((r_0,v_0)\). Let
              \[
                  r_0v_0+I^nV
              \]
              be a neighbourhood of \(r_0v_0\). Suppose
              \[
                  r-r_0\in I^n,\qquad v-v_0\in I^nV.
              \]
              Then
              \[
                  rv-r_0v_0=(r-r_0)v+r_0(v-v_0).
              \]
              Since \(r-r_0\in I^n\), we have
              \[
                  (r-r_0)v\in I^nV.
              \]
              Also, because \(I^nV\) is an \(R\)-submodule of \(V\),
              \[
                  r_0(v-v_0)\in I^nV.
              \]
              Therefore
              \[
                  rv-r_0v_0\in I^nV.
              \]
              Hence
              \[
                  rv\in r_0v_0+I^nV.
              \]
              Thus multiplication \(R\times V\to V\) is continuous.
    
        \item Let \(f\in \operatorname{Hom}_R(V,W)\). We first observe that
              \[
                  f(I^nV)\subseteq I^nW
              \]
              for every \(n\geq 1\). Indeed, if
              \[
                  x=\sum_{j=1}^m a_jv_j\in I^nV,
                  \qquad a_j\in I^n,\ v_j\in V,
              \]
              then
              \[
                  f(x)
                  =
                  f\left(\sum_{j=1}^m a_jv_j\right)
                  =
                  \sum_{j=1}^m a_jf(v_j)
                  \in I^nW.
              \]
              Now let \(v_0\in V\), and let
              \[
                  f(v_0)+I^nW
              \]
              be a neighbourhood of \(f(v_0)\). If
              \[
                  v-v_0\in I^nV,
              \]
              then
              \[
                  f(v)-f(v_0)=f(v-v_0)\in I^nW.
              \]
              Hence
              \[
                  f(v)\in f(v_0)+I^nW.
              \]
              Therefore \(f\) is continuous.
    \end{enumerate}
    \end{proof}
    
    
    \begin{exercise}
    Let \(R\) be a commutative ring, and let \(A\) and \(B\) be \(R\)-algebras.
    Prove the following algebra isomorphisms:
    \begin{enumerate}
        \item
        \[
            M_n(A)\otimes_R B \cong M_n(A\otimes_R B).
        \]
        \item
        \[
            M_n(A)\otimes_R M_m(B) \cong M_{nm}(A\otimes_R B).
        \]
    \end{enumerate}
    \end{exercise}
    
    \begin{proof}
    \begin{enumerate}
        \item Define a map
              \[
                  \Phi:M_n(A)\otimes_R B\longrightarrow M_n(A\otimes_R B)
              \]
              on elementary tensors by
              \[
                  \Phi\bigl((a_{ij})\otimes b\bigr)
                  =
                  (a_{ij}\otimes b).
              \]
              This is well-defined and \(R\)-linear by the universal property of the
              tensor product.
    
              We construct its inverse. Every element of \(M_n(A\otimes_R B)\) can be
              written as a matrix whose \((i,j)\)-entry is a finite sum of elementary
              tensors. Define
              \[
                  \Psi:M_n(A\otimes_R B)\longrightarrow M_n(A)\otimes_R B
              \]
              by
              \[
                  \Psi\bigl((x_{ij})\bigr)
                  =
                  \sum_{i,j}\sum_\ell
                  a_{ij\ell}E_{ij}\otimes b_{ij\ell},
              \]
              whenever
              \[
                  x_{ij}=\sum_\ell a_{ij\ell}\otimes b_{ij\ell}.
              \]
              This is well-defined because tensor product distributes over finite
              direct sums. It is immediate that
              \[
                  \Phi\circ \Psi=\operatorname{id},
                  \qquad
                  \Psi\circ \Phi=\operatorname{id}.
              \]
              Hence \(\Phi\) is an \(R\)-module isomorphism.
    
              It remains to check that \(\Phi\) is an algebra homomorphism. Let
              \(X=(x_{ij})\), \(Y=(y_{ij})\in M_n(A)\), and let \(b,b'\in B\). Then
              \[
                  \Phi\bigl((X\otimes b)(Y\otimes b')\bigr)
                  =
                  \Phi(XY\otimes bb')
                  =
                  ( (XY)_{ik}\otimes bb')_{i,k}.
              \]
              On the other hand,
              \[
                  \Phi(X\otimes b)\Phi(Y\otimes b')
                  =
                  (x_{ij}\otimes b)_{i,j}(y_{jk}\otimes b')_{j,k},
              \]
              whose \((i,k)\)-entry is
              \[
                  \sum_j (x_{ij}\otimes b)(y_{jk}\otimes b')
                  =
                  \sum_j x_{ij}y_{jk}\otimes bb'
                  =
                  (XY)_{ik}\otimes bb'.
              \]
              Thus
              \[
                  \Phi\bigl((X\otimes b)(Y\otimes b')\bigr)
                  =
                  \Phi(X\otimes b)\Phi(Y\otimes b').
              \]
              Also,
              \[
                  \Phi(I_n\otimes 1_B)=I_n.
              \]
              Therefore
              \[
                  M_n(A)\otimes_R B \cong M_n(A\otimes_R B)
              \]
              as \(R\)-algebras.
    
        \item Let \(E_{ij}\) denote the standard matrix units in \(M_n(A)\), and let
              \(F_{kl}\) denote the standard matrix units in \(M_m(B)\). We identify
              the rows and columns of \(M_{nm}(A\otimes_R B)\) with pairs
              \[
                  (i,k),\qquad 1\leq i\leq n,\quad 1\leq k\leq m.
              \]
              Let \(G_{(i,k),(j,l)}\) be the corresponding standard matrix unit in
              \(M_{nm}(A\otimes_R B)\).
    
              Define
              \[
                  \Theta:M_n(A)\otimes_R M_m(B)
                  \longrightarrow
                  M_{nm}(A\otimes_R B)
              \]
              by
              \[
                  \Theta\bigl((aE_{ij})\otimes (bF_{kl})\bigr)
                  =
                  (a\otimes b)G_{(i,k),(j,l)}.
              \]
              Since
              \[
                  M_n(A)=\bigoplus_{i,j} AE_{ij},
                  \qquad
                  M_m(B)=\bigoplus_{k,l} BF_{kl},
              \]
              and tensor product distributes over finite direct sums, \(\Theta\) is an
              \(R\)-module isomorphism.
    
              We check multiplicativity on elementary generators. We have
              \[
                  (aE_{ij})(a'E_{i'j'})
                  =
                  \delta_{j,i'}aa'E_{ij'},
              \]
              and
              \[
                  (bF_{kl})(b'F_{k'l'})
                  =
                  \delta_{l,k'}bb'F_{kl'}.
              \]
              Therefore
              \[
              \begin{aligned}
                  &\Theta\left(
                  \bigl((aE_{ij})\otimes (bF_{kl})\bigr)
                  \bigl((a'E_{i'j'})\otimes (b'F_{k'l'})\bigr)
                  \right) \\
                  &=
                  \delta_{j,i'}\delta_{l,k'}
                  (aa'\otimes bb')G_{(i,k),(j',l')}.
              \end{aligned}
              \]
              On the other hand,
              \[
              \begin{aligned}
                  &\Theta\bigl((aE_{ij})\otimes (bF_{kl})\bigr)
                  \Theta\bigl((a'E_{i'j'})\otimes (b'F_{k'l'})\bigr) \\
                  &=
                  (a\otimes b)G_{(i,k),(j,l)}
                  (a'\otimes b')G_{(i',k'),(j',l')} \\
                  &=
                  \delta_{j,i'}\delta_{l,k'}
                  (aa'\otimes bb')G_{(i,k),(j',l')}.
              \end{aligned}
              \]
              Hence \(\Theta\) preserves multiplication. It also sends the identity to
              the identity. Therefore
              \[
                  M_n(A)\otimes_R M_m(B)
                  \cong
                  M_{nm}(A\otimes_R B)
              \]
              as \(R\)-algebras.
    \end{enumerate}
    \end{proof}